\title{2018 年高考数学 全国III卷-理科}
\date{}
\maketitle

\section*{一、单项选择题}
本题共12小题，每小题5分，共60分。在每小题给出的四个选项中，只有一项是符合题目要求的。

\begin{question}[points=5]
已知集合$A=\{x|x-1\geq 0\}$，$B=\{0,1,2\}$，则$A\cap B=$（ ）
\begin{choices}
  \item $\{0\}$
  \item $\{1\}$
  \item $\{1,2\}$
  \item $\{0,1,2\}$
\end{choices}
\begin{solution}
\textbf{答案：}C

\textbf{分析：}求解不等式化简集合$A$，再由交集的运算性质得答案．

\textbf{详解：}$\because A=\{x|x-1\geq 0\}=\{x|x\geq 1\}$，$B=\{0,1,2\}$，$\therefore A\cap B=\{x|x\geq 1\}\cap\{0,1,2\}=\{1,2\}$．
\end{solution}
\end{question}

\begin{question}[points=5]
$(1+i)(2-i)=$（ ）
\begin{choices}
  \item $-3-i$
  \item $-3+i$
  \item $3-i$
  \item $3+i$
\end{choices}
\begin{solution}
\textbf{答案：}D

\textbf{分析：}直接利用复数代数形式的乘除运算化简得答案．

\textbf{详解：}$(1+i)(2-i)=3+i$．
\end{solution}
\end{question}

\begin{question}[points=5]
中国古建筑借助榫卯将木构件连接起来．构件的凸出部分叫榫头，凹进部分叫卯眼，图中木构件右边的小长方体是榫头．若如图摆放的木构件与某一带卯眼的木构件咬合成长方体，则咬合时带卯眼的木构件的俯视图可以是（ ）
\begin{choices}
  \item A
  \item B
  \item C
  \item D
\end{choices}
\begin{solution}
\textbf{答案：}A

\textbf{分析：}直接利用空间几何体的三视图的画法，判断选项的正误即可．

\textbf{详解：}由题意可知，如图摆放的木构件与某一带卯眼的木构件咬合成长方体，小的长方体，是榫头，从图形看出，轮廓是长方形，内含一个长方形，并且一条边重合，另外3边是虚线，所以木构件的俯视图是A．
\end{solution}
\end{question}

\begin{question}[points=5]
若$\sin\alpha=\frac{1}{3}$，则$\cos2\alpha=$（ ）
\begin{choices}
  \item $\frac{8}{9}$
  \item $\frac{7}{9}$
  \item $-\frac{7}{9}$
  \item $-\frac{8}{9}$
\end{choices}
\begin{solution}
\textbf{答案：}B

\textbf{分析：}$\cos2\alpha=1-2\sin^2\alpha$，由此能求出结果．

\textbf{详解：}$\because \sin\alpha=\frac{1}{3}$，$\therefore \cos2\alpha=1-2\sin^2\alpha=1-2\times\frac{1}{9}=\frac{7}{9}$．
\end{solution}
\end{question}

\begin{question}[points=5]
$(x^2+\frac{2}{x})^5$的展开式中$x^4$的系数为（ ）
\begin{choices}
  \item $10$
  \item $20$
  \item $40$
  \item $80$
\end{choices}
\begin{solution}
\textbf{答案：}C

\textbf{分析：}由二项式定理得展开式的通项，令指数为4求得$r$，再求系数．

\textbf{详解：}由二项式定理得$(x^2+\frac{2}{x})^5$的展开式的通项为：$T_{r+1}=C_5^r (x^2)^{5-r}(\frac{2}{x})^r=2^r C_5^r x^{10-3r}$，由$10-3r=4$，解得$r=2$，$\therefore$系数为$2^2 C_5^2=4\times10=40$．
\end{solution}
\end{question}

\begin{question}[points=5]
直线$x+y+2=0$分别与$x$轴，$y$轴交于$A$，$B$两点，点$P$在圆$(x-2)^2+y^2=2$上，则$\triangle ABP$面积的取值范围是（ ）
\begin{choices}
  \item $[2,6]$
  \item $[4,8]$
  \item $[\sqrt{2},3\sqrt{2}]$
  \item $[2\sqrt{2},3\sqrt{2}]$
\end{choices}
\begin{solution}
\textbf{答案：}A

\textbf{分析：}求出$A$、$B$两点坐标及$|AB|$，设$P$点坐标，利用点到直线距离公式求得距离范围，从而确定面积范围．

\textbf{详解：}由已知$A(-2,0)$，$B(0,-2)$，$|AB|=2\sqrt{2}$．设$P(2+\sqrt{2}\cos\theta,\sqrt{2}\sin\theta)$，则点$P$到直线$x+y+2=0$的距离$d=\frac{|2+\sqrt{2}\cos\theta+\sqrt{2}\sin\theta+2|}{\sqrt{2}}=\frac{|4+2\sin(\theta+\frac{\pi}{4})|}{\sqrt{2}}$，$\because \sin(\theta+\frac{\pi}{4})\in[-1,1]$，$\therefore d\in[\sqrt{2},3\sqrt{2}]$，$\therefore \triangle ABP$面积$\frac{1}{2}\times2\sqrt{2}\times d\in[2,6]$．
\end{solution}
\end{question}

\begin{question}[points=5]
函数$y=-x^4+x^2+2$的图象大致为（ ）

\begin{center}
  \includegraphics[width=.38\linewidth]{question-7-1.jpg}
\end{center}
\begin{choices}
  \item A
  \item B
  \item C
  \item D
\end{choices}
\begin{solution}
\textbf{答案：}D

\textbf{分析：}根据函数图象的特点，求函数的导数利用函数的单调性进行判断．

\textbf{详解：}函数过定点$(0,2)$，排除$A$、$B$．$f'(x)=-4x^3+2x=-2x(2x^2-1)$，令$f'(x)>0$得$x<-\frac{\sqrt{2}}{2}$或$0<x<\frac{\sqrt{2}}{2}$，此时函数递增；令$f'(x)<0$得$x>\frac{\sqrt{2}}{2}$或$-\frac{\sqrt{2}}{2}<x<0$，此时函数递减，排除$C$，故选$D$．
\end{solution}
\end{question}

\begin{question}[points=5]
某群体中的每位成员使用移动支付的概率都为$p$，各成员的支付方式相互独立．设$X$为该群体的10位成员中使用移动支付的人数，$DX=2.4$，$P(X=4)<P(X=6)$，则$p=$（ ）
\begin{choices}
  \item $0.7$
  \item $0.6$
  \item $0.4$
  \item $0.3$
\end{choices}
\begin{solution}
\textbf{答案：}B

\textbf{分析：}利用已知条件，转化为二项分布，利用方差和概率大小关系求解．

\textbf{详解：}由题意$X\sim B(10,p)$，$DX=10p(1-p)=2.4$，解得$p=0.6$或$p=0.4$．由$P(X=4)<P(X=6)$，即$C_{10}^4 p^4(1-p)^6 < C_{10}^6 p^6(1-p)^4$，化简得$(1-p)^2 < p^2$，即$1-p < p$，故$p>0.5$，所以$p=0.6$．
\end{solution}
\end{question}

\begin{question}[points=5]
$\triangle ABC$的内角$A$，$B$，$C$的对边分别为$a$，$b$，$c$．若$\triangle ABC$的面积为$\frac{a^2+b^2-c^2}{4}$，则$C=$（ ）
\begin{choices}
  \item $\frac{\pi}{2}$
  \item $\frac{\pi}{3}$
  \item $\frac{\pi}{4}$
  \item $\frac{\pi}{6}$
\end{choices}
\begin{solution}
\textbf{答案：}C

\textbf{分析：}利用三角形面积公式和余弦定理求解．

\textbf{详解：}$S_{\triangle ABC}=\frac{1}{2}ab\sin C = \frac{a^2+b^2-c^2}{4} = \frac{2ab\cos C}{4} = \frac{ab\cos C}{2}$，所以$\sin C = \cos C$，又$0<C<\pi$，故$C=\frac{\pi}{4}$．
\end{solution}
\end{question}

\begin{question}[points=5]
设$A$，$B$，$C$，$D$是同一个半径为4的球的球面上四点，$\triangle ABC$为等边三角形且面积为$9\sqrt{3}$，则三棱锥$D-ABC$体积的最大值为（ ）
\begin{choices}
  \item $12\sqrt{3}$
  \item $18\sqrt{3}$
  \item $24\sqrt{3}$
  \item $54\sqrt{3}$
\end{choices}
\begin{solution}
\textbf{答案：}B

\textbf{分析：}求出等边三角形边长，确定球心到三角形外心的距离，当$D$在球心与三角形外心连线的延长线上时体积最大．

\textbf{详解：}设等边三角形$ABC$的边长为$a$，则$\frac{\sqrt{3}}{4}a^2=9\sqrt{3}$，解得$a=6$．$\triangle ABC$的外接圆半径$r=\frac{a}{\sqrt{3}}=2\sqrt{3}$，球心$O$到$\triangle ABC$外心$O'$的距离$OO'=\sqrt{R^2-r^2}=\sqrt{16-12}=2$，则$D$到平面$ABC$的最大距离$h=R+OO'=6$，故体积最大值为$\frac{1}{3}\times9\sqrt{3}\times6=18\sqrt{3}$．
\end{solution}
\end{question}

\begin{question}[points=5]
设$F_1$，$F_2$是双曲线$C:\frac{x^2}{a^2}-\frac{y^2}{b^2}=1$（$a>0$，$b>0$）的左、右焦点，$O$是坐标原点．过$F_2$作$C$的一条渐近线的垂线，垂足为$P$，若$|PF_1|=\sqrt{6}|OP|$，则$C$的离心率为（ ）
\begin{choices}
  \item $\sqrt{5}$
  \item $2$
  \item $\sqrt{3}$
  \item $\sqrt{2}$
\end{choices}
\begin{solution}
\textbf{答案：}C

\textbf{分析：}利用点到渐近线的距离求出$|PF_2|$，再求出$|OP|$，在三角形$F_1PF_2$中利用余弦定理建立关系．

\textbf{详解：}双曲线一条渐近线方程为$y=\frac{b}{a}x$，则$F_2(c,0)$到渐近线的距离$d=\frac{bc}{\sqrt{a^2+b^2}}=b$，即$|PF_2|=b$，$|OP|=\sqrt{c^2-b^2}=a$，$\cos\angle PF_2O=\frac{a}{c}$．由$|PF_1|=\sqrt{6}a$，在$\triangle F_1PF_2$中由余弦定理：$|PF_1|^2=|PF_2|^2+|F_1F_2|^2-2|PF_2|\cdot|F_1F_2|\cos\angle PF_2O$，即$6a^2=b^2+4c^2-2b\cdot2c\cdot\frac{a}{c}=b^2+4c^2-4ab$，又$b^2=c^2-a^2$，代入得$6a^2=c^2-a^2+4c^2-4a\sqrt{c^2-a^2}$，化简得$3a^2=c^2$，故$e=\frac{c}{a}=\sqrt{3}$．
\end{solution}
\end{question}

\begin{question}[points=5]
设$a=\log_{0.2}0.3$，$b=\log_2 0.3$，则（ ）
\begin{choices}
  \item $a+b<ab<0$
  \item $ab<a+b<0$
  \item $a+b<0<ab$
  \item $ab<0<a+b$
\end{choices}
\begin{solution}
\textbf{答案：}B

\textbf{分析：}直接利用对数的运算性质化简即可得答案．

\textbf{详解：}$a=\log_{0.2}0.3=\frac{\lg0.3}{\lg0.2}$，$b=\log_2 0.3=\frac{\lg0.3}{\lg2}$，$\because \lg0.3<0$，$\lg0.2<0$，$\lg2>0$，$\therefore a>0$，$b<0$，$ab<0$．又$\frac{1}{a}+\frac{1}{b}=\frac{\lg0.2}{\lg0.3}+\frac{\lg2}{\lg0.3}=\frac{\lg(0.2\times2)}{\lg0.3}=\frac{\lg0.4}{\lg0.3}$，$\because \lg0.4>\lg0.3$，$\therefore \frac{\lg0.4}{\lg0.3}<1$，$\frac{1}{a}+\frac{1}{b}<1$，即$\frac{a+b}{ab}<1$，$\because ab<0$，$\therefore a+b>ab$，故$ab<a+b<0$．
\end{solution}
\end{question}

\section*{二、填空题}
本题共4小题，每小题5分，共20分。

\begin{question}[points=5]
已知向量$\vec{a}=(1,2)$，$\vec{b}=(2,-2)$，$\vec{c}=(1,\lambda)$．若$\vec{c}\parallel(2\vec{a}+\vec{b})$，则$\lambda=$．
\fillin[{$\frac{1}{2}$}]
\begin{solution}
\textbf{答案：}$\frac{1}{2}$

\textbf{分析：}利用向量坐标运算法则求出$2\vec{a}+\vec{b}$，再由向量平行的性质能求出$\lambda$．

\textbf{详解：}$2\vec{a}+\vec{b}=2(1,2)+(2,-2)=(4,2)$，$\because \vec{c}\parallel(2\vec{a}+\vec{b})$，$\therefore$设$\vec{c}=k(4,2)$，即$(1,\lambda)=k(4,2)$，解得$k=\frac{1}{4}$，$\lambda=2k=\frac{1}{2}$．
\end{solution}
\end{question}

\begin{question}[points=5]
曲线$y=(ax+1)e^x$在点$(0,1)$处的切线的斜率为$-2$，则$a=$．
\fillin[{$-3$}]
\begin{solution}
\textbf{答案：}$-3$

\textbf{分析：}求导，利用导数的几何意义列方程求解．

\textbf{详解：}$y'=ae^x+(ax+1)e^x=(ax+a+1)e^x$，则$y'|_{x=0}=a+1=-2$，解得$a=-3$．
\end{solution}
\end{question}

\begin{question}[points=5]
函数$f(x)=\cos(3x+\frac{\pi}{6})$在$[0,\pi]$的零点个数为．
\fillin[{$3$}]
\begin{solution}
\textbf{答案：}$3$

\textbf{分析：}解方程$\cos(3x+\frac{\pi}{6})=0$，求得$x$在$[0,\pi]$内的个数．

\textbf{详解：}由$f(x)=0$得$3x+\frac{\pi}{6}=\frac{\pi}{2}+k\pi$，$k\in\mathbb{Z}$，即$x=\frac{\pi}{9}+\frac{k\pi}{3}$．当$k=0$，$x=\frac{\pi}{9}\in[0,\pi]$；$k=1$，$x=\frac{4\pi}{9}\in[0,\pi]$；$k=2$，$x=\frac{7\pi}{9}\in[0,\pi]$；$k=3$，$x=\frac{10\pi}{9}\notin[0,\pi]$，故零点个数为3．
\end{solution}
\end{question}

\begin{question}[points=5]
已知点$M(-1,1)$和抛物线$C:y^2=4x$，过$C$的焦点且斜率为$k$的直线与$C$交于$A$，$B$两点．若$\angle AMB=90^\circ$，则$k=$．
\fillin[{$2$}]
\begin{solution}
\textbf{答案：}$2$

\textbf{分析：}设直线方程与抛物线联立，利用韦达定理及$\overrightarrow{MA}\cdot\overrightarrow{MB}=0$求解．

\textbf{详解：}焦点$F(1,0)$，设直线$AB:y=k(x-1)$，联立$\begin{cases}y=k(x-1)\\y^2=4x\end{cases}$得$k^2x^2-2(2+k^2)x+k^2=0$，设$A(x_1,y_1)$，$B(x_2,y_2)$，则$x_1+x_2=\frac{2(2+k^2)}{k^2}$，$x_1x_2=1$，$y_1+y_2=k(x_1+x_2-2)=\frac{4}{k}$，$y_1y_2=k^2(x_1-1)(x_2-1)=k^2[x_1x_2-(x_1+x_2)+1]=-4$，$\because \angle AMB=90^\circ$，$\overrightarrow{MA}\cdot\overrightarrow{MB}=0$，即$(x_1+1)(x_2+1)+(y_1-1)(y_2-1)=0$，整理得$x_1x_2+(x_1+x_2)+y_1y_2-(y_1+y_2)+2=0$，代入得$1+\frac{2(2+k^2)}{k^2}-4-\frac{4}{k}+2=0$，解得$k=2$．
\end{solution}
\end{question}

\section*{三、解答题}
共70分。解答应写出文字说明、证明过程或演算步骤。第17~21题为必考题，第22、23题为选考题。

\begin{question}[points=12]
等比数列$\{a_n\}$中，$a_1=1$，$a_5=4a_3$．

（1）求$\{a_n\}$的通项公式；

（2）记$S_n$为$\{a_n\}$的前$n$项和．若$S_m=63$，求$m$．
\begin{solution}

\textbf{分析：}（1）利用等比数列通项公式列出方程，求出公比$q=\pm2$，由此能求出$\{a_n\}$的通项公式．（2）分$q=-2$和$q=2$两种情况，利用等比数列前$n$项和公式求解．

\textbf{详解：}（1）$\because a_1=1$，$a_5=4a_3$，$\therefore 1\cdot q^4=4\times(1\cdot q^2)$，解得$q=\pm2$．当$q=2$时，$a_n=2^{n-1}$；当$q=-2$时，$a_n=(-2)^{n-1}$．（2）当$q=-2$时，$S_n=\frac{1-(-2)^n}{1-(-2)}=\frac{1-(-2)^n}{3}$，令$S_m=63$得$\frac{1-(-2)^m}{3}=63$，即$(-2)^m=-188$，无正整数解；当$q=2$时，$S_n=2^n-1$，令$S_m=2^m-1=63$，解得$m=6$．
\end{solution}
\end{question}

\begin{question}[points=12]
某工厂为提高生产效率，开展技术创新活动，提出了完成某项生产任务的两种新的生产方式．为比较两种生产方式的效率，选取40名工人，将他们随机分成两组，每组20人．第一组工人用第一种生产方式，第二组工人用第二种生产方式．根据工人完成生产任务的工作时间（单位：min）绘制了如下茎叶图：



（1）根据茎叶图判断哪种生产方式的效率更高？并说明理由；

（2）求40名工人完成生产任务所需时间的中位数$m$，并将完成生产任务所需时间超过$m$和不超过$m$的工人数填入下面的列联表：



| | 超过$m$ | 不超过$m$ |

| --- | --- | --- |

| 第一种生产方式 | | |

| 第二种生产方式 | | |



（3）根据（2）中的列联表，能否有99%的把握认为两种生产方式的效率有差异？



附：$K^2=\frac{n(ad-bc)^2}{(a+b)(c+d)(a+c)(b+d)}$，



$P(K^2\geq k)$ | 0.050 | 0.010 | 0.001

--- | --- | --- | ---

$k$ | 3.841 | 6.635 | 10.828

\begin{center}
  \includegraphics[width=.38\linewidth]{question-18-1.jpg}
\end{center}
\begin{solution}

\textbf{分析：}（1）根据茎叶图中的数据判断第二种生产方式的工作时间较少些，效率更高；（2）根据茎叶图中的数据计算它们的中位数，再填写列联表；（3）列联表中的数据计算观测值，对照临界值得出结论．

\textbf{详解：}（1）第二种生产方式效率更高，因为第二种生产方式的工作时间主要集中在65~85之间，而第一种集中在72~92之间，第二种整体工作时间更短．
（2）中位数$m=\frac{79+81}{2}=80$，列联表如下：

| | 超过$m$ | 不超过$m$ | 总计 |
| --- | --- | --- | --- |
| 第一种生产方式 | 15 | 5 | 20 |
| 第二种生产方式 | 5 | 15 | 20 |
| 总计 | 20 | 20 | 40 |

（3）$K^2=\frac{40\times(15\times15-5\times5)^2}{20\times20\times20\times20}=10>6.635$，所以有99%的把握认为两种生产方式的效率有差异．
\end{solution}
\end{question}

\begin{question}[points=12]
如图，边长为2的正方形$ABCD$所在的平面与半圆弧$\\overset{\\frown}{CD}$所在平面垂直，$M$是$\\overset{\\frown}{CD}$上异于$C$，$D$的点．

（1）证明：平面$AMD\perp$平面$BMC$；

（2）当三棱锥$M-ABC$体积最大时，求面$MAB$与面$MCD$所成二面角的正弦值．

\begin{center}
  \includegraphics[width=.38\linewidth]{question-19-1.jpg}
\end{center}
\begin{solution}

\textbf{分析：}（1）根据面面垂直的判定定理证明$MC\perp$平面$ADM$即可．（2）根据三棱锥的体积最大，确定$M$的位置，建立空间直角坐标系，求出点的坐标，利用向量法进行求解．

\textbf{详解：}（1）证明：在半圆中，$DM\perp MC$，$\because$正方形$ABCD$所在的平面与半圆弧$\\overset{\\frown}{CD}$所在平面垂直，$\therefore AD\perp$平面$DCM$，则$AD\perp MC$，$\because AD\cap DM=D$，$\therefore MC\perp$平面$ADM$，又$MC\subset$平面$MBC$，$\therefore$平面$AMD\perp$平面$BMC$．
（2）$\because \triangle ABC$的面积为定值，要使三棱锥$M-ABC$体积最大，则三棱锥的高最大，此时$M$为圆弧$\\overset{\\frown}{CD}$的中点．以$CD$中点$O$为原点，$OC$方向为$x$轴，$OD$方向为$y$轴，过$O$垂直平面$ABCD$向上为$z$轴建立空间直角坐标系，则$A(2,-1,0)$，$B(2,1,0)$，$M(0,0,1)$，平面$MCD$的法向量$\vec{n}_1=(1,0,0)$，$\overrightarrow{AB}=(0,2,0)$，$\overrightarrow{AM}=(-2,1,1)$，设平面$MAB$的法向量$\vec{n}_2=(x,y,z)$，由$\vec{n}_2\cdot\overrightarrow{AB}=2y=0$，$\vec{n}_2\cdot\overrightarrow{AM}=-2x+y+z=0$，令$x=1$，得$y=0$，$z=2$，即$\vec{n}_2=(1,0,2)$，$\cos\langle\vec{n}_1,\vec{n}_2\rangle=\frac{1}{\sqrt{1}\times\sqrt{5}}=\frac{\sqrt{5}}{5}$，故二面角的正弦值为$\sqrt{1-(\frac{\sqrt{5}}{5})^2}=\frac{2\sqrt{5}}{5}$．
\end{solution}
\end{question}

\begin{question}[points=12]
已知斜率为$k$的直线$l$与椭圆$C:\frac{x^2}{4}+\frac{y^2}{3}=1$交于$A$，$B$两点，线段$AB$的中点为$M(1,m)$（$m>0$）．

（1）证明：$k<-\frac{1}{2}$；

（2）设$F$为$C$的右焦点，$P$为$C$上一点，且$\overrightarrow{FP}+\overrightarrow{FA}+\overrightarrow{FB}=0$．证明：$|\overrightarrow{FA}|$，$|\overrightarrow{FP}|$，$|\overrightarrow{FB}|$成等差数列，并求该数列的公差．
\begin{solution}

\textbf{分析：}（1）设$A(x_1,y_1)$，$B(x_2,y_2)$，利用点差法得到$k$与$m$的关系，再由$M$在椭圆内得$m$的范围，从而证明$k<-\frac{1}{2}$．（2）利用向量关系求出$P$点坐标，结合椭圆焦半径公式证明等差数列，并求公差．

\textbf{详解：}（1）设$A(x_1,y_1)$，$B(x_2,y_2)$，则$x_1+x_2=2$，$y_1+y_2=2m$，代入椭圆方程相减得$\frac{x_1^2-x_2^2}{4}+\frac{y_1^2-y_2^2}{3}=0$，即$\frac{(x_1+x_2)(x_1-x_2)}{4}+\frac{(y_1+y_2)(y_1-y_2)}{3}=0$，所以$\frac{2(x_1-x_2)}{4}+\frac{2m(y_1-y_2)}{3}=0$，解得$k=\frac{y_1-y_2}{x_1-x_2}=-\frac{3}{4m}$．由$M(1,m)$在椭圆内得$\frac{1}{4}+\frac{m^2}{3}<1$，解得$0<m<\frac{3}{2}$，故$k=-\frac{3}{4m}<-\frac{1}{2}$．
（2）设$P(x_3,y_3)$，由$\overrightarrow{FP}+\overrightarrow{FA}+\overrightarrow{FB}=0$得$(x_3-1,y_3)+(x_1-1,y_1)+(x_2-1,y_2)=(0,0)$，所以$x_3=1$，$y_3=-(y_1+y_2)=-2m$，代入椭圆得$\frac{1}{4}+\frac{4m^2}{3}=1$，解得$m=\frac{3}{4}$（负舍），故$k=-1$．椭圆右焦点$F(1,0)$，离心率$e=\frac{1}{2}$，焦半径$|FA|=a-ex_1=2-\frac{1}{2}x_1$，$|FB|=2-\frac{1}{2}x_2$，$|FP|=2-\frac{1}{2}x_3=2-\frac{1}{2}=\frac{3}{2}$，$|FA|+|FB|=4-\frac{1}{2}(x_1+x_2)=4-1=3=2\times\frac{3}{2}=2|FP|$，所以$|FA|$，$|FP|$，$|FB|$成等差数列．设直线$AB$方程为$y=-(x-1)$，联立椭圆得$7x^2-8x-8=0$，则$x_1+x_2=\frac{8}{7}$，$x_1x_2=-\frac{8}{7}$，$|x_1-x_2|=\sqrt{(x_1+x_2)^2-4x_1x_2}=\sqrt{\frac{64}{49}+\frac{32}{7}}=\frac{12\sqrt{2}}{7}$，公差$d=\frac{1}{2}|(|FA|-|FB|)|=\frac{1}{2}|(2-\frac{1}{2}x_1)-(2-\frac{1}{2}x_2)|=\frac{1}{4}|x_1-x_2|=\frac{3\sqrt{2}}{7}$，所以公差为$\pm\frac{3\sqrt{2}}{7}$．
\end{solution}
\end{question}

\begin{question}[points=12]
已知函数$f(x)=(2+x+ax^2)\ln(1+x)-2x$．

（1）若$a=0$，证明：当$-1<x<0$时，$f(x)<0$；当$x>0$时，$f(x)>0$；

（2）若$x=0$是$f(x)$的极大值点，求$a$．
\begin{solution}

\textbf{分析：}（1）对函数$f(x)$两次求导数，分别判断$f'(x)$和$f(x)$的单调性，结合$f(0)=0$即可得出结论；（2）令$h(x)$为$f'(x)$的分子，令$h''(0)$计算$a$，讨论$a$的范围，得出$f(x)$的单调性，从而得出$a$的值．

\textbf{详解：}（1）当$a=0$时，$f(x)=(2+x)\ln(1+x)-2x$，$f'(x)=\ln(1+x)+\frac{2+x}{1+x}-2=\ln(1+x)-\frac{x}{1+x}$，$f''(x)=\frac{1}{1+x}-\frac{1}{(1+x)^2}=\frac{x}{(1+x)^2}$，当$-1<x<0$时，$f''(x)<0$，$f'(x)$递减；当$x>0$时，$f''(x)>0$，$f'(x)$递增，所以$f'(x)\geq f'(0)=0$，故$f(x)$在$(-1,+\infty)$上单调递增，又$f(0)=0$，所以当$-1<x<0$时$f(x)<0$，当$x>0$时$f(x)>0$．
（2）$f'(x)=(1+2ax)\ln(1+x)+\frac{2+x+ax^2}{1+x}-2=\frac{(1+2ax)(1+x)\ln(1+x)+ax^2-x}{1+x}$，令$h(x)=(1+2ax)(1+x)\ln(1+x)+ax^2-x$，则$h'(x)=4ax+(4ax+2a+1)\ln(1+x)$，$h''(x)=8a+4a\ln(1+x)+\frac{4ax+2a+1}{1+x}$．由题意$x=0$是极大值点，则$h'(0)=0$，$h''(0)\leq0$，$h''(0)=8a+2a+1=10a+1$，令$10a+1=0$得$a=-\frac{1}{10}$，下面验证：当$a=-\frac{1}{10}$时，$h''(x)=-\frac{4}{5}-\frac{2}{5}\ln(1+x)+\frac{-\frac{2}{5}x-\frac{1}{5}}{1+x}$，$h''(0)=0$，$h'''(x)=-\frac{2}{5(1+x)}+\frac{2}{5(1+x)^2}<0$，故$h''(x)$递减，当$x<0$时$h''(x)>0$，$h'(x)$递增，$h'(x)<h'(0)=0$，$h(x)$递减，$h(x)>h(0)=0$，$f'(x)>0$；当$x>0$时$h''(x)<0$，$h'(x)$递减，$h'(x)<h'(0)=0$，$h(x)$递减，$h(x)<h(0)=0$，$f'(x)<0$，所以$x=0$是极大值点，符合．若$a>-\frac{1}{10}$时，$h''(0)>0$，存在$x>0$使$h'(x)>0$，与极大值矛盾；若$a<-\frac{1}{10}$时，$h''(0)<0$，存在$x<0$使$h'(x)>0$，也与极大值矛盾．综上$a=-\frac{1}{10}$．
\end{solution}
\end{question}

\begin{question}[points=10]
在平面直角坐标系$xOy$中，$\odot O$的参数方程为$\begin{cases}x=\cos\theta\\ y=\sin\theta\end{cases}$（$\theta$为参数），过点$(0,-\sqrt{2})$且倾斜角为$\alpha$的直线$l$与$\odot O$交于$A$，$B$两点．

（1）求$\alpha$的取值范围；

（2）求$AB$中点$P$的轨迹的参数方程．
\begin{solution}

\textbf{分析：}（1）将圆的参数方程化为普通方程，结合直线与圆相交的条件，利用圆心到直线的距离小于半径求解．（2）设直线的参数方程，与圆联立，利用韦达定理和中点坐标公式得到参数方程．

\textbf{详解：}（1）$\odot O$的普通方程为$x^2+y^2=1$，直线$l$的方程为$y=\tan\alpha\cdot x-\sqrt{2}$，圆心$(0,0)$到直线距离$d=\frac{|\sqrt{2}|}{\sqrt{1+\tan^2\alpha}}=\sqrt{2}\cos\alpha<1$，即$\cos\alpha<\frac{\sqrt{2}}{2}$，由于$\alpha\in[0,\pi)$，故$\alpha\in(\frac{\pi}{4},\frac{3\pi}{4})$．
（2）设直线$l$的参数方程为$\begin{cases}x=t\cos\alpha\\ y=-\sqrt{2}+t\sin\alpha\end{cases}$（$t$为参数），代入圆方程得$t^2-2\sqrt{2}t\sin\alpha+1=0$，设$A$、$B$对应参数$t_1$、$t_2$，则$t_1+t_2=2\sqrt{2}\sin\alpha$，所以$P$对应参数$t_P=\frac{t_1+t_2}{2}=\sqrt{2}\sin\alpha$，于是$P$的轨迹参数方程为$\begin{cases}x=\sqrt{2}\sin\alpha\cos\alpha\\ y=-\sqrt{2}+\sqrt{2}\sin^2\alpha\end{cases}$，即$\begin{cases}x=\frac{\sqrt{2}}{2}\sin2\alpha\\ y=-\frac{\sqrt{2}}{2}\cos2\alpha-\frac{\sqrt{2}}{2}\end{cases}$（$\alpha$为参数，$\alpha\in(\frac{\pi}{4},\frac{3\pi}{4})$）．
\end{solution}
\end{question}

\begin{question}[points=10]
设函数$f(x)=|2x+1|+|x-1|$．

（1）画出$y=f(x)$的图象；

（2）当$x\in[0,+\infty)$时，$f(x)\leq ax+b$，求$a+b$的最小值．

\begin{center}
  \includegraphics[width=.38\linewidth]{question-23-1.jpg}
\end{center}
\begin{solution}

\textbf{分析：}（1）利用分段函数的性质将函数表示为分段函数形式进行作图．（2）将不等式恒成立转化为图象关系进行求解．

\textbf{详解：}（1）$f(x)=\begin{cases} -3x, & x\leq -\frac{1}{2} \\ x+2, & -\frac{1}{2}<x<1 \\ 3x, & x\geq 1 \end{cases}$，图象略．
（2）当$x\in[0,+\infty)$时，$f(x)\leq ax+b$恒成立，即$f(x)$的图象在直线$y=ax+b$下方，$x=0$时$f(0)=2\leq b$，$x>0$时，由图象知$a$至少等于$f(x)$在$[0,+\infty)$上各段斜率的最大值3，故$a\geq3$，$b\geq2$，所以$a+b\geq5$，当$a=3$，$b=2$时取等，故$a+b$的最小值为5．
\end{solution}
\end{question}

